\section{Displacement orders and preservation of depth}

Write $v$ for the order on $k[[X]]$, with $v(0)=+\infty$.
A normalized $k$-algebra automorphism $g$ sends
$X$ to $X+\delta$, where $v(\delta)\ge2$.
The following coefficient arguments supply the cancellation details in
\cite{Camina1997}, Lemmas~4--5. Automorphisms act on the left here;
this changes no displacement orders from the paper's right-action convention.

\begin{lemma}[Displacement of a monomial]
\label{lem:normalized-monomial-displacement}
\source{camina1997:page-0006,camina1997:page-0007}
\lean{PowerSeries.order_pow_sub_pow_of_linear_coeff_one}
\leanok
Let $k$ be a field, let $y\in k[[X]]$ have constant coefficient zero
and coefficient of $X$ equal to one, and suppose
$v(y-X)=n+1<\infty$. If $m$ is a natural number whose image in $k$
is nonzero, then
\[
 v(y^m-X^m)=m+n.
\]
\end{lemma}

\begin{proof}
The hypothesis on $m$ implies $m\ge1$. Write $y=Xu$, where
$u(0)=1$. The finite geometric identity gives
\[
 y^m-X^m=X^{m-1}(y-X)(1+u+\cdots+u^{m-1}).
\]
The last factor has constant coefficient $m\ne0$, hence has order zero.
The order of a product of nonzero power series is the sum of their orders.
Consequently the displayed product has order
$(m-1)+(n+1)+0=m+n$.
\end{proof}

\begin{lemma}[Displacement at a leading exponent]
\label{lem:normalized-leading-displacement}
\uses{lem:normalized-monomial-displacement}
\source{camina1997:page-0006,camina1997:page-0007}
\lean{PowerSeries.order_algEquiv_sub_of_order}
\leanok
Let $g$ be a normalized automorphism of $k[[X]]$ with
$v(g(X)-X)=n+1<\infty$. If $v(f)=m<\infty$ and the image
of $m$ in $k$ is nonzero, then
\[
 v(g(f)-f)=m+n.
\]
\end{lemma}

\begin{proof}
Write $f=X^m u$, where $u(0)\ne0$. Every power-series automorphism
preserves constant coefficients, so $g(u)$ has order zero.
Congruence on $X$ modulo $X^{n+1}$ implies the same congruence on
all power series. To see this, it holds for each monomial by
factoring $g(X)^j-X^j$ by $g(X)-X$, and therefore for polynomials;
for a power series, truncate after degree $n$, since both the original
series and its image differ from these polynomials by multiples of
$X^{n+1}$. Thus $v(g(u)-u)\ge n+1$.
Now
\[
 g(f)-f=(g(X)^m-X^m)g(u)+X^m(g(u)-u).
\]
The first summand has order $m+n$ by the preceding lemma, while the
second has order at least $m+n+1$. Their leading terms cannot cancel,
which proves the assertion.
\end{proof}

\begin{lemma}[Displacement of a Frobenius split]
\label{lem:frobenius-split-displacement}
\uses{lem:normalized-leading-displacement}
\source{camina1997:page-0006,camina1997:page-0007}
\lean{PowerSeries.displacement_orders_of_frobenius_split, PowerSeries.le_order_displacement_of_frobenius_split, PowerSeries.order_displacement_of_frobenius_split}
\leanok
Let $p$ be prime and $k$ a field of characteristic $p$. Let $g$ be a normalized
$k$-algebra automorphism of $k[[X]]$, and suppose
$v(g(X)-X)=n+1<\infty$. Let $A,B\in k[[X]]$ satisfy
$v(A)=1$ and $v(B)=u<\infty$, where the image of $u$ in $k$
is nonzero. For $T=A^p+B$,
\[
 v\bigl((g(A)-A)^p\bigr)=p(n+1),\qquad
 v(g(B)-B)=u+n,
\]
and
\[
 v(g(T)-T)\ge\min\{p(n+1),u+n\}.
\]
If $p(n+1)\ne u+n$, the last inequality is an equality.
\end{lemma}

\begin{proof}
The leading-displacement lemma applied first to $A$ and then to
$B$ gives $v(g(A)-A)=n+1$ and $v(g(B)-B)=u+n$.
Multiplicativity of order gives the stated order of the $p$th
power. As $g$ respects addition and powers, and Frobenius is
additive in characteristic $p$,
\[
 g(T)-T=(g(A)-A)^p+(g(B)-B).
\]
The order of a sum is at least the minimum of the two orders.
When these orders differ, the unique term of lower order has
a nonzero leading coefficient which the other summand cannot
cancel. The sum then has precisely that lower order.
\end{proof}

\begin{lemma}[A fixed Frobenius split determines its leading exponent]
\label{lem:frobenius-split-fixed-break}
\uses{lem:frobenius-split-displacement}
\source{camina1997:page-0006}
\lean{PowerSeries.first_exponent_eq_of_fixed_frobenius_split}
\leanok
Under the hypotheses of Lemma~\ref{lem:frobenius-split-displacement},
if $g(A^p+B)=A^p+B$, then $u=n(p-1)+p$.
\end{lemma}

\begin{proof}
If $p(n+1)\ne u+n$, the exact-order conclusion would make
$g(A^p+B)-(A^p+B)$ have finite order, contrary to its being zero.
Hence $p(n+1)=u+n$, and subtracting $n$ gives the asserted formula.
\end{proof}

\begin{lemma}[Depth equality for a Frobenius split]
\label{lem:frobenius-split-depth-equality}
\uses{lem:frobenius-split-displacement}
\source{camina1997:page-0007}
\lean{PowerSeries.depth_eq_of_frobenius_split}
\leanok
Under the hypotheses of Lemma~\ref{lem:frobenius-split-displacement},
suppose there are natural numbers $r,d$ such that
\[
 u=r(p-1)+p,\qquad d<r,\qquad
 v(g(A^p+B)-(A^p+B))=p(d+1).
\]
Then $n=d$.
\end{lemma}

\begin{proof}
If $n\ge r$, then
\[
 p(n+1)\ge p(r+1)>p(d+1),\qquad
 u+n=r(p-1)+p+n\ge p(r+1)>p(d+1).
\]
The lower bound of Lemma~\ref{lem:frobenius-split-displacement}
would contradict the prescribed displacement order. Thus $n<r$.
Since $p>1$, we obtain
\[
 p(n+1)<r(p-1)+p+n=u+n.
\]
The exact-order conclusion of that lemma now gives
$p(d+1)=p(n+1)$. Cancelling the positive factor $p$ gives $d=n$.
\end{proof}

\begin{lemma}[The Frobenius decomposition of a degree-$p$ parameter]
\label{lem:degree-p-parameter-frobenius-decomposition}
\source{camina1997:page-0006,camina1997:page-0007}
\lean{PowerSeries.eq_frobeniusRootPart_pow_add, PowerSeries.order_frobeniusRootPart_of_order_eq_prime, PowerSeries.order_frobeniusRemainder, PowerSeries.exists_frobeniusRemainder_order, PowerSeries.exists_frobeniusRemainder_order_gt_prime}
\leanok
Let $k=\mathbf F_p$ and let
\[
 T=\sum_{i\ge p}a_iX^i,\qquad a_p\ne0.
\]
Suppose $u$ is the least index not divisible by $p$ for which
$a_u\ne0$. Then there are $A,B\in k[[X]]$ such that
\[
 T=A^p+B,\qquad v(A)=1,\qquad v(B)=u,
\]
and the image of $u$ in $k$ is nonzero. If $T$ is not a $p$th
power in $k[[X]]$, such an index $u$ exists and satisfies $u>p$.
\end{lemma}

\begin{proof}
Define
\[
 A=\sum_{j\ge1}a_{pj}X^j,\qquad
 B=\sum_{\substack{i\ge p\\p\nmid i}}a_iX^i.
\]
In characteristic $p$, the $p$th power of a power series has
coefficient $a^p$ in degree $pj$ when the original coefficient
in degree $j$ is $a$, and zero coefficients in other degrees.
Indeed the identity holds for finite truncations by additivity
of Frobenius; every fixed coefficient of the original series and
its $p$th power is determined by a sufficiently long finite
truncation. Since $a^p=a$ in $\mathbf F_p$, the coefficients
of $A^p$ are exactly those of $T$ in degrees divisible by $p$.
The complementary coefficients are precisely those of $B$,
so $T=A^p+B$. The constant term of $A$ vanishes and its
coefficient of $X$ is $a_p\ne0$, giving $v(A)=1$.
The definition of $u$ makes all coefficients of $B$ below $u$
zero and its coefficient of $X^u$ nonzero, giving $v(B)=u$.
Finally $p\nmid u$ says exactly that the image of $u$ in $k$
is nonzero.

For the final assertion, if $B=0$ the identity $T=A^p+B$ would
make $T$ a $p$th power. Thus $B\ne0$, so its nonzero coefficients
have a least index $u$. The defining coefficients of $B$ imply
$p\nmid u$ and that $a_u\ne0$ is the first nonzero coefficient
of $T$ at such an index. Since $v(T)=p$, there are no nonzero
coefficients below $p$, and $p\mid p$ excludes $u=p$.
Hence $u>p$.
\end{proof}

\begin{lemma}[The two possible leading displacement orders]
\label{lem:degree-p-parameter-displacement}
\uses{lem:degree-p-parameter-frobenius-decomposition,lem:frobenius-split-displacement}
\source{camina1997:page-0006,camina1997:page-0007}
\notready
Let $k=\mathbf F_p$ and let
\[
 T=\sum_{i\ge p}a_iX^i,\qquad a_p\ne0.
\]
Suppose $u$ is the least index not divisible by $p$ for which $a_u\ne0$.
Let $g$ be normalized with $v(g(X)-X)=n+1<\infty$.
Then
\[
 v(g(T)-T)\ge\min\{p(n+1),u+n\}.
\]
If $p(n+1)\ne u+n$, equality holds.
\end{lemma}

\begin{proof}
Lemma~\ref{lem:degree-p-parameter-frobenius-decomposition}
gives $T=A^p+B$ with $v(A)=1$, $v(B)=u$, and $u$ nonzero
in $k$. Apply Lemma~\ref{lem:frobenius-split-displacement}
to obtain both assertions.
\end{proof}

\begin{lemma}[The first exponent detects the degree-p ramification break]
\label{lem:camina-four-depth-equation}
\uses{lem:degree-p-parameter-frobenius-decomposition,lem:frobenius-split-fixed-break}
\source{camina1997:page-0006}
\notready
Under the preceding hypotheses, if $g(T)=T$, then
\[
 u=n(p-1)+p.
\]
\end{lemma}

\begin{proof}
Take the decomposition $T=A^p+B$ supplied by
Lemma~\ref{lem:degree-p-parameter-frobenius-decomposition}.
Its orders satisfy all the hypotheses of
Lemma~\ref{lem:frobenius-split-fixed-break}, and $g(T)=T$ is
its fixed-point hypothesis. That lemma gives the formula.
\end{proof}

\begin{lemma}[Depth below a new ramification break is unchanged]
\label{lem:camina-five-depth-preservation}
\uses{lem:degree-p-parameter-frobenius-decomposition,lem:frobenius-split-depth-equality}
\source{camina1997:page-0007}
\notready
Under the degree-p parameter hypotheses, suppose that for a natural number $d$
\[
 g(T)=T+\sum_{j\ge d+1}\beta_jT^j,\qquad \beta_{d+1}\ne0,
\]
and $u=r(p-1)+p$ for a natural number $r>d$.
If $v(g(X)-X)=n+1<\infty$, then $n=d$.
\end{lemma}

\begin{proof}
Since $v(T)=p$, write the given displacement as
$T^{d+1}V(T)$ with $V(Z)=\beta_{d+1}+\beta_{d+2}Z+\cdots$.
The substitution converges because $v(T)>0$, and $V(T)$
has constant coefficient $\beta_{d+1}\ne0$, hence order zero.
Thus $v(g(T)-T)=p(d+1)$.
Take the decomposition $T=A^p+B$ of
Lemma~\ref{lem:degree-p-parameter-frobenius-decomposition}.
Lemma~\ref{lem:frobenius-split-depth-equality} applies:
its remaining hypotheses are $u=r(p-1)+p$ and $d<r$.
It gives $n=d$ as required.
\end{proof}
